% Fresh derivation, 2026-10-22. No missing historical IW witness is imported.
\section{The inverse-$b$ uniform C1 source: a complete local conductor}
\label{sec:c1-inverse-b-local-new}

\paragraph{Statement and certificate ring.}
Let $K_c,K_t$ be independent real closed extensions of $\mathbb R$,
with bounded rings $B_c,B_t$ and infinitesimal ideals
$\mathfrak m_c,\mathfrak m_t$. The certificate ring is
$\mathcal T=K_c\otimes_{\mathbb R}K_t$, with bounded coefficient ring
$\mathcal B=B_c\otimes_{\mathbb R}B_t$ and zero-residue ideal
$\mathfrak n=\mathfrak m_c\otimes B_t+B_c\otimes\mathfrak m_t$.
We use the finite strict-residue Gram lemma in this ring.
All radicals and inverses below are pure source, pure testing, or
products of separately specified such factors.

Use the actual uniform native C1 stationary source, with
$\lambda=n=0$, $d>0$, and the original nonnegative event coefficients.
Normalize $d=1$ by a positive pure-source scale. Suppose
\[
 k,f,g>0\text{ are infinitesimal},\qquad k/g\in\mathfrak m_c,
 \qquad r=1-af^2,\quad h=1-tg^2,
\]
\[
 R=1+k^2r,\quad T_0=1+k^2h,\quad
 0<a<T_0,\quad 0<t<R,\quad \beta\geq0.                 \tag{IW1}
\]
All original strict gaps remain imposed. Then the complete original
parent is tensor-SOS whenever the independent positive testing angles
$u,v$ are infinitesimal, for every real or complex testing state.
The proof below works in the rotated chart $w=1,z>0$ with $z$ arbitrary;
in particular it includes the previously missing original
$W$-unit, $Z$-infinitesimal chart. It uses no bound on $g/f$ and no
comparison between any testing and source infinitesimals.
The exact sign, zero-axis and state reductions are given at the end.

\subsection{The actual C1 source and bounded source arrays}

For clarity, the six weights used here come from the following
unchanged native formulas. Put $x_n=k^2$, $\eta=f^2$, $\zeta=g^2$,
$y_n=x_n\eta\zeta$, $K=1+x_n$, $G=h+\zeta$, and
\[
 M=R-t,\quad J=T_0-a,\quad E_n=1-x_n\eta h,
 \quad Q_n=1-x_n\zeta r,\quad \theta=t\zeta.
\]
Define
\[
 P_d=1+2x_nh+rx_n^2G+x_n\zeta
       +\eta x_n(K\zeta+x_nh\theta),
\]
\[
 P_4=T_0E_n+x_n(1-r)G,\qquad
 P_\alpha=Rh(1+y_n)+\zeta M E_n.
\]
The actual relative coefficients, now with $d=1$, are
\[
\begin{array}{ll}
 m=\dfrac{Kx_n(r+\eta)h(1+y_n)}{2M\zeta P_d},&
 C=\dfrac{K(r+\eta)E_n}{2(1-r)P_d},\\[4pt]
 A=\dfrac{\eta P_4}{2rP_d},&
 D+\tfrac12=\dfrac{K(1+y_n)(r+\eta)}{2J\eta P_d},\\[4pt]
 \beta+\tfrac12=\dfrac{Kx_n(r+\eta)G}{2\theta P_d},&
 \alpha=\dfrac{x_n^2\zeta(r+\eta)P_\alpha}{2Q_nP_d}.
\end{array}                                                     \tag{IW2}
\]
No auxiliary C2 stationary source is used.
Since $a,t$ are bounded and $k,f,g$ are infinitesimal,
$r,h,P_d,R,T_0$ have residue one. The inequality
$\beta+\tfrac12\geq\tfrac12$ in (IW2) gives
$t=O(k^2/g^2)$. Thus $t$ and $p=t/R$ are infinitesimal, and
$M=R(1-p)$ is a positive unit. The two independent openings
$a$ and $J=T_0-a$ are not capped.

Introduce the exact compact coordinates
\[
 j=\frac{T_0}{g^2R},\quad B=\frac{k^2g^2R}{T_0},\quad
 T=\frac{f^2T_0}{R},\quad p=\frac tR,\quad H=\frac a{T_0},
 \qquad x=jB,\quad y=TB,\quad q_z=jT.                    \tag{IW3}
\]
Here $j$ is unbounded, $0<H<1$, and
\[
 T,x,B,p,\sigma:=1/j,\delta:=j^2B\in\mathfrak m_c,
 \qquad \delta\sim k^2/g^2,\qquad p=O(\delta).
                                                               \tag{IW4}
\]
The quantity $q_z$ can have any positive size. The exact inverse is
\[
\begin{gathered}
 R=\frac{1+x}{1+HTx},\quad T_0=\frac{1+x}{1+Bp},\quad
 r=\frac{1-HT}{1+HTx},\quad h=\frac{j-p}{j(1+Bp)},\\
 f^2=\frac{T(1+Bp)}{1+HTx},\quad
 g^2=\frac{1+HTx}{j(1+Bp)},\quad a=HT_0,\quad t=pR.
\end{gathered}                                                   \tag{IW5}
\]
In particular both strict gaps are retained exactly.

Here is a short rational dictionary sufficient to check every source
estimate. Write $D_+=D+\tfrac12$, $b_+=\beta+\tfrac12$,
$\mu=m(1-p)$, and set
\[
\begin{split}
 P={}&1+B(j+1-p)+BT(1+Hj)\\
 &+B^2T(Hj+jp-p^2+p)+B^2HT^2j(1+Bp),\\
 L={}&1+T(1-H+Bp),\qquad E=1+B(HT+p-1),\\
 F_\alpha={}&j+1-2p+BHTj(1-p)+BTp(j-p).
\end{split}
\]
Substitution of (IW5) in (IW2) gives exactly
\[
\begin{split}
 \mu&=\frac{Bj(j-p)(1+y)(1+Bp)L}{2(1+x)P},\\
 C&=\frac{(1+Bp)L[1+BTp-BTj(1-H)]}{2HT(1+x)P},\\
 A&=\frac{T(1+Bp)[1+BT(2Hj-Hp+H-j+p)+B^2H^2T^2j]}
              {2(1-HT)P},\\
 D_+&=\frac{(1+y)(1+Bp)(1+HTx)L}{2T(1-H)(1+x)P},\\
 b_+&=\frac{Bj(1+Bp)L(j+1-p+HTx)}{2p(1+x)P},\\
 \alpha&=\frac{B^2LF_\alpha}{2EP}.
\end{split}                                                     \tag{IW6}
\]
The denominators $P,E,1+x,1-HT$ are positive bounded units.
After $j=1/\sigma$, $B=\delta\sigma^2$, each of the following
normalized arrays is a rational function with polynomial numerator
and denominator in the bounded variables $\sigma,\delta,T,p,H$;
its denominator has a nonzero real constant at
$\sigma=\delta=T=p=0$, independently of $H$:
\[
 \frac{2\mu}{\delta},\quad\frac{2b_+p}{\delta},\quad
 2TCH,\quad2TD_+(1-H),\quad\frac{2A}{T},\quad
 \frac{2\alpha}{B^2j}\quad\hbox{all have residue }1.
                                                               \tag{IW7}
\]
This finite polynomial criterion proves boundedness and the stated
residues at every relative rate and at both $H$ openings. It also
shows $\beta p\leq b_+p=O(\delta)$, including $\beta=0$.
A useful exact cancellation is
\[
                  K_d:=CH-D_+(1-H)=-\frac{b_+p}{j}.    \tag{IW8}
\]

\subsection{The entire parent and both state normals}

Put $D_u=1+u^2$, $D_v=1+v^2$. The quarter-parent is
\[
\begin{split}
 \mathcal P={}&m(vW-X+uvY-uZ)^2+CD_v(uW-Y)^2+AD_uD_vY^2\\
 &+D_+(uW+uvX-Y-vZ)^2
 +(\alpha+\tfrac12)D_u(vY-Z)^2\\
 &+b_+D_v(X+uZ)^2-\tfrac12D_uD_vZ^2.
\end{split}                                                     \tag{IW9}
\]
On the rotated chart
$w=(W+uY)/D_u=1$, $z=(Z-uX)/D_u>0$, use the two unrestricted normals
\[
 N=\frac{X+uZ-pD_uv}{D_u},\qquad
 M_0=\frac{uW-Y-HD_uvz}{D_u}.
\]
Their exact inverse is
\[
 X=N+pv-uz,\quad Z=z+uN+puv,\quad
 W=1+u(M_0+Hvz),\quad Y=u-M_0-Hvz.                     \tag{IW10}
\]
Only the pure-testing unit $D_u$ is inverted. Retain the positive
packet $mpD_u^2N^2$, and put $G=(\mathcal P-mpD_u^2N^2)/D_u$.
Direct substitution, before any state specialization, gives
\[
\begin{split}
 G={}&\mu D_u[N^2-2vN+(1-p)v^2]
      +CD_uD_v(M_0+Hvz)^2+A D_v(u-M_0-Hvz)^2\\
 &+D_+D_u(M_0-(1-H)vz)^2\\
 &+(\alpha+\tfrac12)
       [u(1-p)v-(1+Hv^2)z-uN-vM_0]^2\\
 &+b_+D_uD_v(N+pv)^2-\tfrac12D_v(z+uN+puv)^2.
\end{split}                                                     \tag{IW11}
\]
In particular $G$ is exactly quadratic in the two normals:
\[
 G=F+2L_NN+2L_MM_0+H_NN^2+H_MM_0^2
                         +2(\alpha+\tfrac12)uvNM_0,
                                                               \tag{IW12}
\]
where $F=G|_{N=M_0=0}$ and
\[
\begin{split}
 h_0&=\mu+b_+\geq\tfrac12,\\
 H_N&=h_0D_u+b_+v^2+\beta u^2v^2+\alpha u^2,\\
 H_M&=D_u(CD_v+D_+)+AD_v+(\alpha+\tfrac12)v^2.
\end{split}                                                     \tag{IW13}
\]
The potentially unbounded $\beta$ is fully retained in (IW11)--(IW13).
There is no omitted higher-normal polynomial.

\subsection{An exact separated scalar circuit}

Define the rows
\[
 U=u/\sqrt x,\quad V=v/\sqrt y,\quad Z_*=z/\sqrt{q_z},
 \quad Q_*=U-VZ_*,\quad B_*=V-UZ_*,\quad R_*=Z_*-UV,
\]
\[
 \Phi=U^2+V^2+Z_*^2+U^2V^2Z_*^2-4UVZ_*.
\]
None of $U,V,Z_*$ is assumed bounded. On the positive chart,
\[
 \Phi=(U-V)^2+Z_*^2(1-UV)^2+2UV(1-Z_*)^2             \tag{IW14}
\]
is a legal separated SOS; for example
$\sqrt{UV}=(xy)^{-1/4}\sqrt{uv}$.

For an explicit finite definition of the actual masses, put
$K_0=CH^2+D_+(1-H)^2$ and
\[
\begin{array}{ll}
 M_a=Ax,&M_b=[\mu(1-p)+b_+p^2]y,\\
 M_c=\alpha q_z,&
 M_{bc}=[K_0+AH^2+2H\alpha+H-\tfrac12]yq_z,\\
 M_{b^2}=b_+p^2y^2,&
 M_{b^2c}=H^2(A+C+\alpha+\tfrac12)y^2q_z,\\
 M_{ab^2c}=CH^2xy^2q_z,&M_{ab^2}=\beta p^2xy^2.
\end{array}
\]
Set
\[
\begin{gathered}
 f_c=M_b,\quad h_c=M_{ab^2},\quad
 k_c=(M_a-M_b+M_{ab^2c}-M_{bc}+h_c)/2,\\
 r_c=M_c-M_b-k_c,\quad q_c=M_a-f_c,\quad
 g_c=M_{ab^2c},\quad d_c=M_{b^2c}-k_c,
 \quad e_c=b_+p^2y^2-g_c.
\end{gathered}
\]
The complete scalar from (IW11) satisfies the exact identity
\[
\begin{split}
 F={}&q_cQ_*^2+r_cR_*^2+d_cV^2Q_*^2+e_cV^2B_*^2\\
 &+f_c\Phi+g_cV^2\Phi+k_cZ_*^2(V^2-1)^2+h_cV^2R_*^2.
\end{split}                                                     \tag{IW15}
\]
The last term is the extra nonnegative square supplied by the actual
$\beta$ event. The checker verifies all scalar monomials of the
original parent and the four remaining source coefficient identities;
(IW15) is not merely an identity at the contact.

The same finite bounded-array criterion as (IW7), with the displayed
mass definitions, gives
\[
 \frac{2f_c}{Tx^2}\sim1,\quad
 \frac{2q_c}{Tx}\sim1,\quad
 \frac{2r_c}{TjB^2}\sim1,\quad
 \frac{k_c}{T^2B^2j}\in\mathfrak m_c,\quad
 \frac{g_c}{f_c}\in\mathfrak m_c.                     \tag{IW16}
\]
Here $\sim1$ means residue one. In particular the three principal
masses $f_c,q_c,r_c$ are positive, and
$|k_c|/r_c$, $|k_c|/(q_cy)$ are infinitesimal.
All normalized mass denominators remain units after
$j=1/\sigma$, $B=\delta\sigma^2$; the companion finite-array
checks test their residues independently of $H$.

For the signed terms use the exact identities
\[
\begin{split}
 (1+U^2)\Phi&=R_*^2+B_*^2+U^2Q_*^2+U^2(1-UVZ_*)^2,\\
 (1+V^2)\Phi&=Q_*^2+R_*^2+V^2B_*^2+V^2(1-UVZ_*)^2,\\
 Z_*(V^2-1)&=-VQ_*-R_*.
\end{split}                                                     \tag{IW17}
\]
The $e_c,g_c$ combination in (IW15) is exactly
\[
 b_+p^2y^2V^2B_*^2-g_c\Phi
       +g_c[Q_*^2+R_*^2+V^2(1-UVZ_*)^2].
\]
A $k_c$ of either sign costs at most
$3|k_c|V^2Q_*^2+2|k_c|R_*^2$, including its occurrence in $d_c$.
The corresponding relative costs are
$3|k_c|v^2/(q_cy)$ and $2|k_c|/r_c$, both in $\mathfrak n$.
Every inequality here abbreviates a finite SOS identity, using
$2A^2+2B^2-(A+B)^2=(A-B)^2$ and nonnegative pure-source factors.
Consequently the scalar supplies strict fixed real fractions of
\[
                       q_cQ_*^2,\qquad r_cR_*^2,
                       \qquad f_c\Phi.                \tag{IW18}
\]
All other nonnegative packets remain available as whole squares.

\subsection{Finite budgets for the complete normal rows}

Call a row free if any fixed positive real multiple of its square can
be subtracted at infinitesimal cost from (IW18), by a finite separated
SOS identity. This definition allows a bounded normal coupling to be
absorbed by a sufficiently large fixed real helper pivot.
Put
\[
 q=vz-Tu,\quad b=v-uz/j,\quad R_z=z-uv/B,\quad
 R_u=v(u^2-x),\quad R_v=u(v^2-y),
\]
\[
 J_u=u(u^2-x),\qquad J_v=v(v^2-y),\qquad
 \widehat q=q/\sqrt T.
\]
Since $Q_*=-q/(T\sqrt x)$, the first mass in (IW18) has a positive
unit coefficient on $\widehat q^2$.

First, $R_u=\sqrt y\,x(-B_*-UR_*)$. Applying the first identity
of (IW17) to $B_*^2$, and the actual $r_c$ pivot to $U^2R_*^2$,
shows that $R_u$ is free. The three relative costs are bounded by
constant multiples of
\[
                         B,\qquad u^2/j,\qquad u^2,
                                                               \tag{IW19}
\]
all infinitesimal. Second,
$R_v=\sqrt x\,y(-Q_*-VR_*)$, so $\sqrt T R_v$ is free;
its two costs are $O(y^2)$ and $O(Tv^2)$.

The same-angle rows have exact decompositions
\[
\begin{split}
 J_v&=\sqrt y\,v(v+\sqrt{y/x}\,u)(V-U)+(y/x)R_u,\\
 J_u&=\sqrt x\,u(u+\sqrt{x/y}\,v)(U-V)+(x/y)R_v.
\end{split}                                                     \tag{IW20}
\]
If $\rho=O(\delta)$ is a nonnegative source factor, then
$\rho J_v$ is free. Indeed $\rho\sqrt{y/f_c}=O(\sqrt\delta)$,
$\sqrt{y/x}=\sqrt{T/j}$ is bounded, and
$\rho y/x$ is infinitesimal. Thus every coefficient after division
by the strict $\sqrt{f_c}(V-U)$ pivot is infinitesimal.
Also $\sqrt T\,TxJ_u$ is free: the two coefficients of its
first summand relative to $\sqrt{f_c}(U-V)$ are bounded multiples of
\[
                  T\sqrt x\,u^2,\qquad\sqrt{T\delta}\,uv,
\]
and its second summand is exactly
$\delta(\sqrt T R_v)$, because $Tx^2/y=j^2B=\delta$.
This last use of $\delta\in\mathfrak m_c$ is the additional
source fact specific to the inverse-$b$ sector.

Finally,
\[
 j\alpha b=-(\alpha/B)R_u-\alpha uR_z.                \tag{IW21}
\]
Here $\alpha/B=O(x)$, and the square of $\alpha uR_z$ costs
$O(x^2u^2)$ relative to $r_cR_*^2$. Hence the entire left side
of (IW21) is free. These identities involve no bound on $z$.

\paragraph{The unrestricted first normal.}
Define the bounded source quantities
\[
 K_z=H(\alpha+\tfrac12)-\tfrac12,\qquad
 K_R=-\mu-\alpha(1-p)-\tfrac12+b_+p+TK_z,
\]
\[
 \rho=b_+p+\beta px\geq0.
\]
Direct differentiation of the entire (IW11), followed by the actual
source contact identity, gives
\[
 L_N=-j\alpha b+(K_R+\beta pv^2)R_u
                   +\rho J_v+K_zuvq.                 \tag{IW22}
\]
For completeness, before imposing the source, the remainder is
$v S_N$, with
\[
 S_N=j\alpha-\mu+b_+p+K_Rx+b_+py+\beta pxy=0
\]
identically under (IW6). Thus no constant or first-normal term is
silently dropped. Equations (IW7) imply
$K_R\sim-\tfrac12$, $\rho=O(\delta)$ and
$\beta p=O(\delta)$. The first three terms in (IW22) are free by
(IW19)--(IW21); the last has coefficient
$K_z\sqrt T\,uv\in\mathfrak n$ on $\widehat q$.
Dividing by $\sqrt{h_0}$ preserves all these bounds because
$h_0\geq\tfrac12$.

\paragraph{The unrestricted second normal.}
Use the exact $K_d$ from (IW8), and put
\[
\begin{split}
 C_q={}&K_dD_u+AH+\alpha+\tfrac12
        +v^2[CHD_u+AH+(\alpha+\tfrac12)H],\\
 c_2={}&T[CH+AH+(\alpha+\tfrac12)H]
                      -A-(\alpha+\tfrac12)(1-p).
\end{split}
\]
The complete second normal has the exact division
\[
 L_M=C_qq+(c_2+TCHu^2)R_v+T(K_d+CHy)J_u.             \tag{IW23}
\]
Before substituting the actual source its only remainder is $uS_M$,
where
\[
\begin{split}
 c_0&=T[K_d+AH+\alpha+\tfrac12]-A,\\
 S_M&=c_0+TK_dx+yc_2+TCHxy=0.
\end{split}
\]
The last equality is an exact identity from (IW6). Furthermore
$c_2$ is bounded (in fact infinitesimal), $TCH$ is a unit, and
$(K_d+CHy)/x$ is bounded with residue $-\tfrac12$.
It follows from the two free rows in (IW20) that
\[
 \sqrt{TH(1-H)}[(c_2+TCHu^2)R_v+T(K_d+CHy)J_u]
\]
is free with bounded coefficients. For the other term,
\[
 \sqrt{TH(1-H)}C_qq
     =T\sqrt{H(1-H)}C_q\widehat q.
\]
Its coefficient is in $\mathfrak n$: multiplying $C_q$ by $T$
clears the only $CH$ pole, and every term then contains either
$T$ or $v^2$. This statement is made on the full expression in
(IW23), before any relative source/testing rate is considered.

\subsection{Strict full Gram and state restoration}

Normalize the hard rows by pure-source factors,
\[
 \widehat N=\sqrt{h_0}\,N,\qquad
 M_0=\sqrt{TH(1-H)}\,\widehat M.
\]
Their full hard matrix is bounded: $b_+/h_0\leq1$,
$\beta/h_0\leq1$, $\alpha/h_0=O(\alpha)$, and the two
source poles of $H_M$ are exactly cleared by $TH(1-H)$.
Equations (IW7), (IW13) give its real residue
\[
                           \operatorname{diag}(1,\tfrac12).
                                                               \tag{IW24}
\]
Its off-diagonal coefficient has the factor
$uv\sqrt{TH(1-H)}/\sqrt{h_0}$ and is infinitesimal.
The $\widehat q$ couplings in (IW22)--(IW23) are infinitesimal.
All remaining complete-normal couplings are bounded coefficients
of free rows. Buy sufficiently large fixed real pivots on the
finitely many free rows, use the elementary real Young square
completion, and apply the strict-residue Gram lemma to the
remaining finite matrix. The strict scalar pivots (IW18) lose only
infinitesimal amounts. This proves $G\in\Sigma\mathcal T^2$.
No state normal has been set to zero in this conclusion.
Restoring $mpD_u^2N^2$, the pure-testing factor $D_u$, and the
positive source scale $d$ proves the assertion for (IW9).

The proof has used $u,v>0$ and $w=1,z>0$, with arbitrary $z$.
The exact sign/state lemma SJ1--SJ10 applies to the same original
core: sign gauges make $u,v$ nonnegative while preserving every
packet; a negative core at positive angles forces the rotated
$w,z$ to have the same nonzero sign. Dividing all state coordinates
by the nonzero pure-testing $w$ gives precisely this chart.
When the core is nonnegative it is itself a pure-testing square,
and every event packet is already positive. Zero angles have the
direct core SOS exits. The zero state is immediate, and a complex
state splits into its real and imaginary quadratic evaluations.
Thus (IW1) is proved for all states at the double-infinitesimal
angles. The existing two-current outer-angle identities give the
separate all-tests join if needed; they are not required as a premise
of the local conductor proved here.

\paragraph{Verification and scope.}
The new module \texttt{iw\_actual\_source.py} prints the finite
actual C1 source dictionary. The independent-formula checks in
\texttt{check\_c1\_inverse\_b\_local.py} compare it with (IW2),
verify the entire original parent and both normals, all scalar
monomials and the exact circuit, the elementary helper identities,
and the finite normalized source arrays. The ordered-field
estimates and finite Gram allocation are the proof above; a checker
receipt alone does not establish them. No lost historical IW
coefficient witness, R21 admission label, formal negative-$\beta$
source, mixed inverse, or additional source-rate restriction is used.
